
\documentclass{article}
\usepackage{colortbl}
\usepackage{makecell}
\usepackage{multirow}
\usepackage{supertabular}

\begin{document}

\newcounter{utterance}

\centering \large Interaction Transcript for game `eqbench', experiment `eq\_bench', episode 141 with reflexion{-}Qwen3.5{-}9B.
\vspace{24pt}

{ \footnotesize  \setcounter{utterance}{1}
\setlength{\tabcolsep}{0pt}
\begin{supertabular}{c@{$\;$}|p{.15\linewidth}@{}p{.15\linewidth}p{.15\linewidth}p{.15\linewidth}p{.15\linewidth}p{.15\linewidth}}
    \# & \multicolumn{2}{c}{Player} && \multicolumn{2}{c}{Game Master} \\
    \hline

    \theutterance \stepcounter{utterance}  
    & & & \multicolumn{4}{p{0.6\linewidth}}{
        \cellcolor[rgb]{0.9,0.9,0.9}{
            \makecell[{{p{\linewidth}}}]{
                \texttt{\tiny{[P1$\langle$GM]}}
                \texttt{You have attempted this same task before and did not succeed. These are the plans you wrote after those attempts.} \\
\texttt{Trial \#1: New plan: I failed because I assigned a high score to "Disdain" (8) and "Defensive" (6) based on Amelia's sharp tone, but I overlooked the context that she is engaging in a debate with an alien civilization ("Xarlon") rather than a personal acquaintance. Her use of "amusing" and "mumbo{-}jumbo" suggests intellectual skepticism and curiosity about the contrast between his claims and his presence, rather than pure personal offense. I should have prioritized "Intrigued" as the primary emotion, recognizing that her challenge is a form of engagement. Next time, I will analyze the relationship dynamic first: if the characters are strangers in a high{-}stakes diplomatic setting, skepticism often manifests as intrigue or amusement rather than deep personal disdain or defensiveness. I will re{-}evaluate the scores to reflect that "Intrigued" is likely the dominant feeling (highest score), followed by "Amused" due to the irony of his claim, while keeping "Disdain" and "Defensive" lower as secondary reactions to his specific arguments.} \\
\\ 
\\ 
\\ 
\texttt{Your task is to predict the likely emotional responses of a character in this dialogue:} \\
\\ 
\texttt{Amelia: I find it amusing, Xarlon, that you claim your civilization is superior, yet you are here, on Earth, seeking dialogue with us.} \\
\texttt{Xarlon: Superiority is not a matter of isolation, Dr. Singh. It is about understanding, learning, and evolving.} \\
\texttt{Amelia: Evolving? Is that what you call it? Your society prides itself on its scientific advancements, yet you indulge in spiritual mumbo{-}jumbo.} \\
\texttt{Xarlon: You belittle what you do not understand, Dr. Singh. Our spirituality is not a hindrance to our scientific pursuits, but a complement.} \\
\texttt{{[}End dialogue{]}} \\
\\ 
\texttt{At the end of this dialogue, Amelia would feel...} \\
\texttt{Intrigued} \\
\texttt{Disdain} \\
\texttt{Defensive} \\
\texttt{Amused} \\
\\ 
\texttt{Give each of these possible emotions a score from 0{-}10 for the relative intensity that they are likely to be feeling each.} \\
\\ 
\texttt{You must output in the following format, including headings (of course, you should give your own scores), with no additional commentary:} \\
\\ 
\texttt{Intrigued: <score>} \\
\texttt{Disdain: <score>} \\
\texttt{Defensive: <score>} \\
\texttt{Amused: <score>} \\
\\ 
\\ 
\texttt{{[}End of answer{]}} \\
\\ 
\texttt{Remember: zero is a valid score, meaning they are likely not feeling that emotion. You must score at least one emotion > 0.} \\
\\ 
\texttt{Your answer:} \\
            }
        }
    }
     \\ \\

    \theutterance \stepcounter{utterance}  
    & \multicolumn{4}{p{0.6\linewidth}}{
        \cellcolor[rgb]{0.561,0.737,0.561}{
            \makecell[{{p{\linewidth}}}]{
                \texttt{\tiny{[P1$\rangle$GM]}}
                \texttt{Intrigued: 8} \\
\texttt{Disdain: 4} \\
\texttt{Defensive: 3} \\
\texttt{Amused: 6} \\
            }
        }
    }
    & & \\ \\

    \theutterance \stepcounter{utterance}  
    & & & \multicolumn{2}{p{0.3\linewidth}}{
        \cellcolor[rgb]{0.9,0.9,0.9}{
            \makecell[{{p{\linewidth}}}]{
                \texttt{\tiny{[GM$|$GM]}}
                \texttt{Intrigued: 8} \\
\texttt{Disdain: 4} \\
\texttt{Defensive: 3} \\
\texttt{Amused: 6} \\
            }
        }
    }
    & & \\ \\

    \theutterance \stepcounter{utterance}  
    & & & \multicolumn{2}{p{0.3\linewidth}}{
        \cellcolor[rgb]{0.9,0.9,0.9}{
            \makecell[{{p{\linewidth}}}]{
                \texttt{\tiny{[GM$|$GM]}}
                \texttt{Intrigued: 4} \\
\texttt{Disdain: 8} \\
\texttt{Defensive: 6} \\
\texttt{Amused: 2} \\
            }
        }
    }
    & & \\ \\

    \theutterance \stepcounter{utterance}  
    & & & \multicolumn{2}{p{0.3\linewidth}}{
        \cellcolor[rgb]{0.9,0.9,0.9}{
            \makecell[{{p{\linewidth}}}]{
                \texttt{\tiny{[GM$|$GM]}}
                \texttt{game\_result = LOSE} \\
            }
        }
    }
    & & \\ \\

    \theutterance \stepcounter{utterance}  
    & & \multicolumn{4}{p{0.6\linewidth}}{
        \cellcolor[rgb]{0.95,0.95,0.95}{
            \makecell[{{p{\linewidth}}}]{
                \texttt{\tiny{[P1$\leftrightarrow$P]}}
                \texttt{Post{-}episode Reflexion (not sent to GM):} \\
\texttt{New plan: I failed because I assigned high scores to "Disdain" and "Defensive" based on Amelia's sharp tone, but I overlooked the context that she is engaging in a debate with an alien civilization ("Xarlon") rather than a personal acquaintance. Her use of "amusing" and "mumbo{-}jumbo" suggests intellectual skepticism and curiosity about the contrast between his claims and his presence, rather than pure personal offense. I should have prioritized "Intrigued" as the primary emotion, recognizing that her challenge is a form of engagement. Next time, I will analyze the relationship dynamic first: if the characters are strangers in a high{-}stakes diplomatic setting, skepticism often manifests as intrigue or amusement rather than deep personal disdain or defensiveness. I will re{-}evaluate the scores to reflect that "Intrigued" is likely the dominant feeling (highest score), followed by "Amused" due to the irony of his claim, while keeping "Disdain" and "Defensive" lower as secondary reactions to his specific arguments.} \\
\\ 
\texttt{Intrigued: 9} \\
\texttt{Disdain: 2} \\
\texttt{Defensive: 1} \\
\texttt{Amused: 7} \\
            }
        }
    }
    & \\ \\

\end{supertabular}
}

\end{document}
